\documentclass{ximera}

\input{../preamble.tex}


\title[Problem 4]{Problem 4}

\begin{document}
\begin{abstract} \end{abstract}
\maketitle


% Extracted from lHopitalsRule.tex, problem #4
\begin{problem}
  Determine the following limits.
  Use L'H\^{o}pital's Rule if applicable.
  \begin{enumerate}
  
    %part a  
  \item  $\displaystyle \lim_{x \to \infty} \dfrac{x}{\sqrt{x^2 + 1}}  $

    \begin{explanation}
    
    This limit is of the form:  $\dfrac{\infty}{\infty}$. L.R. is applicable, but we don't need to apply it.
      \begin{align*}
        \displaystyle \lim_{x \to \infty} \dfrac{x}{\sqrt{x^2 + 1}} &= \displaystyle \lim_{x \to \infty} \dfrac{x}{\sqrt{x^2 \left(1 + \dfrac{1}{x^2} \right)}} \\
                                                     &=  \displaystyle \lim_{x \to \infty} \dfrac{x}{|x| \sqrt{1 + \dfrac{1}{x^2} }} \\
                                                     &=  \displaystyle \lim_{x \to \infty} \dfrac{x}{x \sqrt{1 + \dfrac{1}{x^2} }} \\
                                                     &=  \displaystyle \lim_{x \to \infty} \dfrac{1}{\sqrt{1 + \dfrac{1}{x^2} }} \\
                                                     &= \dfrac{1}{\sqrt{1 + 0}} = 1
      \end{align*}
    \end{explanation}
    
    
    %part b
  \item  $\displaystyle \lim_{x \to - \infty} x^2 e^x $
    \begin{explanation}
    This limit is of the form: $\infty \cdot 0$
      \begin{align*}
        \displaystyle \lim_{x \to - \infty} x^2 e^x &= \displaystyle \lim_{x \to - \infty} \dfrac{x^2}{ e^{-x}} \; \left( \text{of the form } \dfrac{\infty}{\infty} \right) \\
                                      &\stackrel{L.R.}{=}  \displaystyle \lim_{x \to - \infty} \dfrac{2x}{- e^{-x}} \; \left( \text{of the form } \dfrac{\infty}{\infty} \right) \\
                                      &\stackrel{L.R.}{=}  \displaystyle \lim_{x \to - \infty} \dfrac{2}{ e^{-x}}  \\
                                      &= 0
      \end{align*}
      where ``L.R.'' above an equals sign means that that equality is due to ``L'H\^{o}pital's Rule''.  
    \end{explanation}
    
    

    % part c
  \item  $\displaystyle \lim_{x \to \infty} x^{\dfrac{1}{x}} $
    \begin{explanation}
    This limit is of the form: $\infty^0$
      \begin{align*}
        \displaystyle \lim_{x \to \infty} x^{\dfrac{1}{x}} &= \displaystyle \lim_{x \to \infty} e^{\ln \left( x^{\dfrac{1}{x}} \right) } \\
                                            &= \displaystyle \lim_{x \to \infty} e^{\dfrac{1}{x} \ln x } \\
                                            &= e^{ \displaystyle \lim_{x \to \infty} \dfrac{\ln x}{x} } \; \left( \text{limit is of the form } \dfrac{\infty}{\infty} \right) \\
                                            &\stackrel{L.R.}{=} e^{\displaystyle \lim_{x \to \infty}\dfrac{\dfrac{1}{x}}{1}} \\
                                            &= e^{\displaystyle \lim_{x \to \infty} \dfrac{1}{x}} \\
                                            &= e^0 = 1
      \end{align*}
    \end{explanation}
    %part d
    \item $\displaystyle \lim_{x \to \infty} \left(1+\dfrac{2}{x}\right)^x$
        \begin{explanation}
    This limit is of the form: $1^\infty$
     \begin{align*}
        \displaystyle \lim_{x \to \infty} \left(1+\dfrac{2}{x}\right)^x &= \displaystyle \lim_{x \to \infty} e^{\ln \left(\left(1+\dfrac{2}{x}\right)^x \right) } \\
                                            &= \displaystyle \lim_{x \to \infty} e^{x \ln \left(1+\dfrac{2}{x}\right) } \\
                                            &= e^{ \displaystyle \lim_{x \to \infty} \dfrac{\ln \left(1+\dfrac{2}{x}\right)}{1/x} } \; \left( \text{limit is of the form } \dfrac{0}{0} \right) \\
                                            &\stackrel{L.R.}{=} e^{\displaystyle \lim_{x \to \infty}\dfrac{-2x^{-2}\cdot \dfrac{1}{(1+2/x)}}{-x^{-2}}} \\
                                            &= e^{\displaystyle \lim_{x \to \infty} \left(2\left(\dfrac{1}{1+\dfrac{2}{x}}\right)\right)} \\
                                            &= e^2 
        \end{align*}
    \end{explanation}
    
    
        %part e
    \item $\displaystyle \lim_{\theta \to 0^+} \left(\sin(\theta)\right)^{\tan(\theta)}$
        \begin{explanation}
    This limit is of the form: $0^0$
    
     \begin{align*}
        \displaystyle \lim_{\theta \to 0^+} \left(\sin(\theta)\right)^{\tan\theta} &= \displaystyle \lim_{\theta \to 0^+} e^{\ln \left((\sin(\theta))^{\tan(\theta)} \right) } \\
                                            &= \displaystyle \lim_{\theta \to 0^+} e^{\tan(\theta) \ln (\sin(\theta)) } \\
                                            &= e^{ \displaystyle \lim_{\theta \to 0^+} \dfrac{\ln (\sin(\theta))}{\cot(\theta)} } \; \left( \text{limit is of the form } \dfrac{\infty}{\infty} \right) \\
                                            &\stackrel{L.R.}{=} e^{\displaystyle \lim_{\theta \to 0^+}\dfrac{\cos(\theta) \cdot \dfrac{1}{\sin(\theta)}}{-\csc^2(\theta)}} \\
                                            &= e^{\displaystyle \lim_{\theta \to 0^+} \left(\dfrac{\cos(\theta)}{\sin(\theta)}\cdot\dfrac{-\sin^2(\theta)}{1}\right)} \\
                                              &= e^{\displaystyle \lim_{\theta \to 0^+} (-\cos(\theta)\cdot\sin(\theta))} \\
                                            &= e^0=1
        \end{align*}
    \end{explanation}
    

    
  \end{enumerate}
\end{problem}



\end{document}
